% AI-produced research fragment; no priority claim.  Reserved IDs A.12--A.24.
\subsection{The Legendre representation: the exact proved range}

Throughout this subsection a \emph{Selmer section} means a full arithmetic
section whose finite local restrictions are point sections or cuspidal
sections, and whose real restriction has the usual real-component
interpretation.  Write
\[
 \mathcal E/U:\ y^2=x(x-1)(x-t),\qquad
 L_{s,\ell}=s^*T_\ell\mathcal E,\qquad
 V_{s,\ell}=L_{s,\ell}\otimes\Q_\ell.
\]
The maps defining these representations are the actual continuous
composites $G_\Q\xrightarrow{s}\pi_1(U)\to
\operatorname{GL}_2(\Z_\ell)$.  Reduction gives
$\bar\rho_{s,\ell}:G_\Q\to\operatorname{GL}_2(\mathbb F_\ell)$.
The Weil pairing identifies their determinants with the cyclotomic
characters.  The arguments below use the density-one integrality result
AP-D.25 in its Selmer-section range, and finite-support recognition
LM6-B.14--B.15.  These attached results were opened, including their proofs
\cite{LegAP2026,LegLM62026}.  Their passage from Selmer sections to arbitrary
sections requires the separate local-section audit; no such passage is
silently included in the following statements.

The finite-support input has the following precise use here.  At every
finite etale neighbourhood of a Selmer section, take its local centres on
the smooth completion.  They survive finite descent, by pulling back a
finite torsor along the section and taking the resulting arithmetic twist.
Stoll's finite-support theorem puts a density-one tuple supported on a
fixed finite subscheme into that subscheme's rational points.  Applied
through the neighbourhoods, the compatible rational pro-point gives an
interior point section or a cusp decomposition section.  Only finite
places are required in this support argument; the correction concerning
real components is retained \cite[Theorem 8.2 and Proposition 8.5]{LegStoll2007}
\cite[\S3]{LegStollErrata2017}.

\begin{statement}{A.12}{Proved}{Rational absolute irreducibility}
For every noncuspidal Selmer section $s$ of $U/\Q$ and every prime
$\ell$, the representation $V_{s,\ell}$ is absolutely irreducible.
This assertion does not assume finite ramification.
\end{statement}
\begin{proof}
At the coefficient prime, a point restriction is the Tate module of an
actual elliptic curve, hence potentially semistable with Hodge--Tate
weights $0,-1$, in the convention that $\Q_\ell(1)$ has weight $-1$.
At a cusp, after its quadratic splitting twist, it is an extension of
$\Q_\ell$ by $\Q_\ell(1)$.  Every such extension is semistable, including
the split extension and every completed Kummer class; the explicit
period calculation is in \cite[II.4.4, pp.~18--19]{LegBerger2002}.
Undoing the finite twist preserves potential semistability and the weights.

We need a character calculation that also allows a finite coefficient
extension $F/\Q_\ell$.  A continuous character
$\psi:G_\Q\to F^\times$ that is Hodge--Tate at $\ell$ has the form
\[
                        \psi=\eta\chi_\ell^a,
       \qquad a\in\Z,\quad \eta\text{ of finite order}.
\tag{A.12.1}
\]
Indeed, compactness puts its image in $\mathcal O_F^\times$.
Choose an integer $h>0$ whose power map sends this unit group into a
torsion-free pro-$\ell$ group.  Kronecker--Weber identifies
$G_\Q^{\rm ab}=\prod_q\Z_q^\times$.  For $q\ne\ell$, the restriction
of $\psi^h$ kills both the finite part and the open pro-$q$ part of
$\Z_q^\times$.  The direct sum of these factors is dense in their
product.  Thus $\psi^h$ factors through $\Z_\ell^\times$.
Tate's local algebraicity theorem
\cite[III, Appendix A.7, Theorem 3]{LegSerre1989}, applied to the
underlying $\Q_\ell$-representation of the $F$-line, makes the inertia
action algebraic on an open subgroup.  It commutes with multiplication
by $F$, hence its algebraic cocharacter factors through
$\operatorname{Res}_{F/\Q_\ell}\mathbb G_m$.  A cocharacter defined over
$\Q_\ell$ has one common integral exponent at all embeddings of $F$,
because those embeddings form a single Galois orbit.  Consequently
$\eta=\psi\chi_\ell^{-a}$ is trivial on an open subgroup of inertia.
The preceding power argument makes $\eta^h$ finite, hence $\eta$ finite,
proving (A.12.1).

If $V_{s,\ell}$ had a line over $\overline\Q_\ell$, the line would be
defined over a finite coefficient extension.  Subobjects and quotients
of a Hodge--Tate representation are Hodge--Tate with additive weight
multiplicities \cite[Proposition 2.4.4]{LegBrinonConrad2009}.
Equation (A.12.1), the two weights, and the determinant would therefore
give semisimplification $\eta\oplus\eta^{-1}\chi_\ell$.
The primes splitting in the finite extension killing $\eta$ have
positive density.  They meet AP-D.25's density-one set of integral
point restrictions.  At every such prime $p\ne\ell$ the trace is
$p+1$, whereas the good elliptic fibre has trace of absolute value at
most $2\sqrt p$.  The strict inequality $p+1>2\sqrt p$ is a
contradiction.
\end{proof}

\begin{statement}{A.13}{Proved}
{Residual irreducibility and its exact coefficient range}
For a noncuspidal Selmer section, $\bar\rho_{s,\ell}$ is absolutely
irreducible for every $\ell\ge5$.  At $\ell=3$ this fails for an actual
rational point section; at $\ell=2$ the representation is trivial for
every section.
\end{statement}
\begin{proof}
Let $Z_\ell=\mathbb P(\mathcal E[\ell])\to U$, the finite etale
degree-$\ell+1$ cover of cyclic subgroups.  An invariant line is a fixed
point in the finite $G_\Q$-set pulled back along $s$.  Choosing it gives
an actual section lift to the associated open subgroup
$\pi_1(Z_\ell)\subset\pi_1(U)$.  The two geometric monodromies
$\left(\begin{smallmatrix}1&2\\0&1\end{smallmatrix}\right)$ and
$\left(\begin{smallmatrix}1&0\\-2&1\end{smallmatrix}\right)$ generate
$\operatorname{SL}_2(\mathbb F_\ell)$ for odd $\ell$; its transitive
action on $\mathbb P^1(\mathbb F_\ell)$ proves geometric connectedness.

Here is the required morphism to the modular curve.  On the universal
fibre let $P=(0,0)$ and $Q=(1,0)$.  For
$\phi:E\to E'=E/\langle P\rangle$ and its dual
$\phi':E'\to E$, put $H_4=(\phi')^{-1}(\langle Q\rangle)$.
If $\phi'(R)=Q$, then $2R=\phi(Q)\ne0$; hence $H_4$ is cyclic of
order four.  If $C_\ell$ is the universal subgroup on $Z_\ell$, then
$H_4+\phi(C_\ell)$ is cyclic of order $4\ell$.  This defines
\[
             Z_\ell\longrightarrow Y_0(4\ell)
\tag{A.13.1}
\]
over $\Q$.  It is nonconstant, since a degree-two isogeny has finitely
many possible source $j$-invariants for a fixed target, whereas
$j(E_t)$ is nonconstant.  It extends to a finite map of completions.

The lifted local sections remain point or cusp sections: a fixed sheet
over an ordinary point is a rational point of its etale fibre; over a
cusp the decomposition-group intersection selects a rational boundary
point.  Their completed local tuple survives finite descent.  The index
of $\Gamma_0(4\ell)$ is $6(\ell+1)$, its cusp count is six, and its
elliptic-point counts are zero, so its genus is $(\ell-3)/2>0$.
Stoll's Corollary 8.8 therefore places the image tuple in
$X_0(4\ell)(\Q)$ \cite{LegStoll2007}.  Every such rational point is a
cusp.  For the precise rational-isogeny exclusion needed here, Mazur's
table and composite-degree discussion already suffice: the only five
then-unsettled minimal degrees listed on p.~131 are
$169,91,65,39,125$, none of which divides $4\ell$ for prime $\ell\ge5$;
the recorded degrees and their isogeny graphs contain no $4\ell$
\cite[pp.~129--131]{LegMazur1978}.  Thus no unopened composite-isogeny
classification is being substituted at this step.

The finite inverse image of the rational cusp in (A.13.1) is contained
in the boundary of $Z_\ell$.  Finite-support recognition makes the
lifted section cuspidal, hence makes $s$ cuspidal, a contradiction.
This proves irreducibility over $\mathbb F_\ell$.  An irreducible
two-dimensional representation over this field that is not absolutely
irreducible splits over its algebraic closure into two Frobenius-conjugate
characters.  Their values at complex conjugation are equal, since both
belong to $\{1,-1\}\subset\mathbb F_\ell$.  Their determinant there
would be $+1$, contradicting the cyclotomic determinant $-1$.

For $\ell=3$, take $t=32/27$.  The division polynomial is
\[
 \psi_3(x)=3x^4-4(1+t)x^3+6tx^2-t^2.
\]
It vanishes at $x=4/3$, where the curve value is $y^2=16/243$.
The two points with this rational $x$-coordinate are negatives, so their
order-three subgroup is $G_\Q$-stable.  This is a reducible residual
representation of a noncuspidal rational point section.  Finally,
$O,(0,0),(1,0),(t,0)$ identify the finite etale group scheme
$\mathcal E[2]$ with the constant group $(\Z/2)^2_U$.  Pullback along
any section is therefore trivial.  These claims concern the explicit
reduction maps; they do not contradict A.12.
\end{proof}

\begin{statement}{A.14}{Proved}{The repaired bounded-support equivalence}
For noncuspidal Selmer sections of $U/\Q$, the following four properties
define the same class:
\begin{enumerate}[label=(\roman*)]
\item One finite set contains the prime-to-$\ell$ ramification supports
of all $\bar\rho_{s,\ell}$, $\ell\ge5$.
\item $V_{s,\ell}$ is finitely ramified for at least one $\ell\ge5$.
\item The local sections are points $x_p\in U(\Z_p)$ at all but
finitely many primes.
\item One finite set contains the prime-to-$\ell$ ramification
supports of all $V_{s,\ell}$.
\end{enumerate}
Every section in this class has a rational Tate-system realization
$V_{s,\ell}\simeq V_\ell(E)$ for one elliptic curve $E/\Q$ and every
$\ell$.  The statement does not assert that all noncuspidal sections
belong to the class.
\end{statement}
\begin{proof}
First verify the finite-ramification route.  Under (ii), A.12 and its
local calculation give irreducibility, oddness, potential semistability,
and distinct Hodge--Tate weights.  These, \emph{together with finite
ramification}, are the hypotheses of Pan's Theorem 1.0.4, read together
with Conjecture 1.0.1 in the opened author version
\cite{LegPan2022}.  Its extra residual exclusion concerns only
$\ell=3$.  The theorem gives eigenform origin up to a character twist.
That character is a de Rham subobject of the corresponding Hom
representation.  The argument (A.12.1) makes it finite order times a
cyclotomic power.  Matching the weights forces weight two and zero
cyclotomic power; absorb the finite character into the eigenform.
The determinant then gives trivial nebentype.

Outside the finite ramification set, the coefficient prime, 2, and the
modular level, consider a local section.  An interior point gives an
unramified Tate module, hence good reduction by Neron--Ogg--Shafarevich.
For the Legendre formula
\[
             j(t)=256\frac{(t^2-t+1)^3}{t^2(1-t)^2},
\tag{A.14.1}
\]
any odd-prime value outside $U(\Z_p)$ has negative $p$-valuation of
$j$.  This contradicts good reduction.  A local cusp gives
$\delta_p\otimes\left(\begin{smallmatrix}\chi_\ell&*\\0&1
\end{smallmatrix}\right)$ with $\delta_p$ quadratic.  Unramifiedness
of the whole representation forces $\delta_p$ unramified, so its trace
is $\pm(p+1)$.  Modular weight-one purity gives absolute value at most
$2\sqrt p$, another contradiction.  This includes the split cuspidal
extension.  Thus (ii) implies (iii).  The elliptic scheme over an
integral Legendre point gives (iii)$\Rightarrow$(iv); (iv) implies (ii)
and, by reduction of the stable lattice, (i).

For completeness, rationality of the Hecke field and realization of the
whole system follow as follows.  On the density-one integral set the
good-prime eigenvalues are integers.  Any conjugate eigenform has the
same traces there.  Chebotarev and continuity give equal characters,
and newform multiplicity one gives equality of the forms
\cite[Theorem 9.4 and the paragraph after Example 9.5]{LegStein2007}.
The coefficient field is $\Q$.  Eichler--Shimura gives an elliptic
curve \cite[Proposition 11.21]{LegMilneMF2017}.  At every other
coefficient prime, equality of the traces on the same density-one set,
A.12, and Brauer--Nesbitt give $V_{s,\ell}\simeq V_\ell(E)$.

It remains to prove (i)$\Rightarrow$(iv), since the direct contact-order
argument in L4 does not cover general cuspidal Kummer parameters.
Enlarge the finite set $S$ in (i) by 2 and 3.  Let $D$ be the density-one
set of integral point restrictions and put
$a_p=p+1-\#E_{x_p}(\mathbb F_p)$ for $p\in D$, $p\ne2$.
Choose coefficient primes $\ell\in D\setminus S$ in the CRT class
\[
 \ell\equiv5\pmod {72},\qquad
 \ell\equiv2\pmod p\quad(p\in S,\ p\ge5).
\tag{A.14.2}
\]
This is a unit residue class, so Dirichlet and the density-one property
give infinitely many such primes.  At $\ell$ the residual representation
comes from the $\ell$-torsion of a good-reduction elliptic curve and is
finite flat; its Serre weight is two.  A.13 gives absolute
irreducibility.  The order
$|\operatorname{GL}_2(\mathbb F_\ell)|=\ell(\ell-1)^2(\ell+1)$
and (A.14.2) bound the wild inertia image by $2^5$ at 2, by 3 at 3,
and trivially at all other primes of $S$.  The local conductor estimate
in Helm--Voloch's Proposition 6 gives
\[
 N_\ell\mid N:=2^{14}3^5\prod_{p\in S,\ p\ge5}p^2.
\tag{A.14.3}
\]
That calculation uses only the residual ramification support and finite
wild inertia image; it does not assume finite ramification of the
original adic representation \cite[Theorem 5 and Proposition 6]{LegHelmVoloch2011}.

Serre modularity therefore supplies weight-two forms of levels dividing
$N$.  Among the finitely many normalized newforms in these spaces one
form $f$, with coefficient field $K$, occurs for infinitely many
distinct residue characteristics $\ell$.  For fixed $p\in D\setminus S$,
the algebraic integer $a_p(f)-a_p$ lies in prime ideals of $K$ over
infinitely many rational primes.  Its norm forces it to be zero.
The same comparison of determinants gives trivial nebentype on $D$,
hence everywhere by Chebotarev.  The preceding conjugation and
multiplicity-one argument again gives $K=\Q$ and an elliptic curve $E$.

Fix any coefficient prime $r$.  The traces of $V_r(E)$ and $V_{s,r}$
agree on $D$ away from a finite set.  This implies equality of their
continuous characters even before finite ramification of $V_{s,r}$ is
known.  For each finite quotient of the two images and every conjugacy
class in that quotient, Chebotarev supplies a positive-density
Frobenius set; removing the complement of $D$ leaves such Frobenius
elements.  Equality on all finite quotients and continuity give equality
on $G_\Q$.  A.12 and Brauer--Nesbitt give
$V_{s,r}\simeq V_r(E)$, proving (iv).
\end{proof}
The unproved assertion in A4 is that the union in (i) is finite for an
arbitrary noncuspidal section.  A.14 proves the equivalence claimed by
L4/L5, with the missing global argument supplied, but no bound on that
union.

\begin{statement}{A.15}{Proved}
{Prime residual supports do not detect arbitrary cusp contacts}
There is a globally cuspidal section for which every prime residual
support, away from its coefficient prime, is contained in $\{2\}$,
while every adic Legendre representation is infinitely ramified.
\end{statement}
\begin{proof}
Enumerate the rational primes as $p_j$ and put
$b=(\ell)_\ell\in\prod_\ell\Z_\ell=\widehat\Z$ and
$a_{p_j}=j!b$.  For each integer $n$, all but finitely many $a_{p_j}$
vanish modulo $n$.  A.10 gives a completed rational multiplicative
class $A$ with these valuation coordinates and positive sign.
Every reduction modulo a prime $\ell$ is zero.  For fixed $\ell$,
however, the $\ell$-adic valuation at every $p_j$ is $j!\ell\ne0$.
Twist the tangent-one splitting of the cusp-zero decomposition extension
by the continuous Kummer cocycle of $A$.
Since $q(t)=t^2/256+O(t^3)$, the Tate extension changes by
$2\kappa_\ell(A)$ after the fixed quadratic splitting twist.
For $p_j\ne2,\ell$ the reference class and that twist are unramified.
Residual inertia is zero, whereas adic tame inertia has exponent
$2j!\ell\ne0$.  This proves the assertion with the actual reduction
map $\widehat\Z\to\mathbb F_\ell$; its simultaneous kernel is
$\prod_\ell\ell\Z_\ell$.
\end{proof}

\subsection{Integral realization, its pairing obstruction, and local rigidity}

\begin{statement}{A.16}{Proved}{An integral Legendre realization}
Every noncuspidal Selmer section in the class of A.14 has
$\lambda_0\in\Q\setminus\{0,1\}$ and an isomorphism
\[
             \prod_\ell L_{s,\ell}\simeq
             \prod_\ell T_\ell(E_{\lambda_0})
\tag{A.16.1}
\]
of continuous $\widehat\Z[G_\Q]$-modules respecting the ordered
level-two structure.  The elliptic curve here is an actual Legendre
curve over $\Q$; the quadratic twist is detected and removed below.
\end{statement}
\begin{proof}
Start with the rational realization $E$ from A.14.  At $\ell\ge5$,
A.13 makes the lattice reduction irreducible.  Any two stable lattices
in its rational representation are homothetic: rescale one into the
other with nonzero reduction, use irreducibility for surjectivity of
that reduction, and then use Nakayama.  Thus the lattices can be
identified integrally at every $\ell\ge5$.
At 2 and 3 rescale the transported prescribed lattice to contain
$T_\ell(E)$.  Its finite quotient is a $G_\Q$-stable subgroup of
$V_\ell(E)/T_\ell(E)\simeq E[\ell^\infty]$.  Quotient $E$ by the
product of these two finite subgroups.  Pullback along this isogeny
identifies the resulting Tate lattices with the prescribed lattices at
2 and 3, and changes none at the other primes.  Call the resulting
curve $E'$.  Its two labeled rational 2-torsion points give a monic
model
\[
                 y^2=(x-e_1)(x-e_2)(x-e_3),\qquad e_i\in\Q,
\]
with $P_1=(e_1,0)$ and $P_2=(e_2,0)$ the ordered pair.

Choose $a^2=e_1-e_2$, $b^2=e_1-e_3$, $c^2=e_2-e_3$, and $i^2=-1$.
The points
\[
 R_1=(e_1+ab,-ab(a+b)),\qquad
 R_2=(e_2+iac,-iac(ia+c))
\]
satisfy $2R_1=P_1$, $2R_2=P_2$, so form a basis of $E'[4]$.
Let $A,B,C,I$ denote the sign characters on $a,b,c,i$ with values in
$\mathbb F_2$.  The chord formula gives
\[
 \sigma R_1-R_1=A P_1+(A+B)P_2,\qquad
 \sigma R_2-R_2=(I+A+C)P_1+(I+A)P_2.
\]
For example, changing $b$ to $-b$ adds $P_2$ to $R_1$;
changing $a$ adds $P_1+P_2$.  Applying the same calculation with
base point $P_2$ and the roots $ia,c$ gives the second formula.
The matrix is consequently
\[
 \begin{pmatrix}1+2A&2(I+A+C)\\2(A+B)&1+2(I+A)\end{pmatrix}
                  \pmod4.
\tag{A.16.2}
\]
The lower-right character is $1+2\kappa_2(e_2-e_1)$, since
$ia$ is a square root of $e_2-e_1$.  Replacing the halves while keeping
their ordered doubles conjugates by $1+2M$; conjugation of $1+2N$ by
$1+2M$ is trivial modulo four.  Thus this character is an invariant of
the labeled integral representation.

For the universal Legendre curve take $e_1=0,e_2=1,e_3=t$.
Then $A=I$, $B=I+\kappa_2(t)$, $C=\kappa_2(1-t)$, and (A.16.2) is
\[
       \begin{pmatrix}\chi_4&2\kappa_2(1-t)\\
                       2\kappa_2(t)&1\end{pmatrix}.
\tag{A.16.3}
\]
Its lower-right entry is identically one.  Pullback by $s$ and the
labeled integral comparison therefore give $e_2-e_1=d^2$ for some
$d\in\Q^\times$.  The actual rational coordinate change
\[
 x=e_1+d^2X,\qquad y=d^3Y,
 \qquad \lambda_0=\frac{e_3-e_1}{e_2-e_1}
\]
takes $E'$ to $Y^2=X(X-1)(X-\lambda_0)$ and respects the ordered
2-torsion.  This proves (A.16.1).
\end{proof}

\begin{statement}{A.17}{Proved}
{The exact geometric-conjugator obstruction}
In the standard Legendre basis the closure of geometric monodromy is
\[
 H=\prod_{\ell\ne2}\operatorname{SL}_2(\Z_\ell)\ \times\
 \left\{\begin{pmatrix}a&b\\c&d\end{pmatrix}\in
 \operatorname{SL}_2(\Z_2): b,c\in2\Z_2,
                     \ a,d\equiv1\pmod4\right\}.
\tag{A.17.1}
\]
For a labeled integral conjugator $g$ in A.16, the class of its Weil
multiplier
\[
                 [\det g]\in
           \widehat\Z^\times/(\widehat\Z^\times)^2
\tag{A.17.2}
\]
is independent of $g$.  There is a conjugator in $H$ exactly when this
class is trivial.  Its triviality for every section in A.14 has not
been proved here.
\end{statement}
\begin{proof}
Write $U(x)=\left(\begin{smallmatrix}1&x\\0&1\end{smallmatrix}\right)$
and $L(y)=\left(\begin{smallmatrix}1&0\\y&1\end{smallmatrix}\right)$.
Geometric monodromy is generated by $U(2),L(-2)$.  Taking their powers
with exponents in $\widehat\Z$ produces $U(2z),L(2z)$ for every
$z\in\widehat\Z$, including idempotents supported at one coefficient
prime.  At an odd prime these generate all of $\operatorname{SL}_2$.
At 2 they lie in the displayed subgroup.  Conversely, Gauss
decomposition reduces a matrix of that subgroup to
$\operatorname{diag}(a,a^{-1})$ with $a\equiv1\pmod4$ and even
upper and lower factors.  The identity
\[
 \operatorname{diag}(a,a^{-1})=
 U(2)L((a-1)/2)U(-2/a)L(-a(a-1)/2)
\]
expresses the diagonal factor with even arguments.  Taking closures
of finite products proves (A.17.1).

By A.12 and Schur's lemma, two integral $G_\Q$-isomorphisms in A.16
differ at each coefficient prime by a scalar in $\Z_\ell^\times$.
Their determinants therefore differ by a square, proving independence
of (A.17.2).  Necessity follows from $H\subset\operatorname{SL}_2$.
For sufficiency, rescale $g$ by a scalar square root of its inverse
determinant.  The result has determinant one and still reduces to the
identity modulo two.  Its two 2-adic diagonal entries are equal modulo
four, because its determinant is one and its off-diagonal entries are
even.  If both are $-1$ modulo four, multiply the 2-component by $-I$.
This changes neither determinant nor conjugacy and puts it in $H$.
Thus there is no additional independent 2-adic sign obstruction.
\end{proof}
Equality of determinants with the cyclotomic character for two
representations says that an intertwiner is a \emph{similitude}; it
does not itself compute its multiplier.  The exact homomorphism at
issue is the determinant map on the scalar torsor of intertwiners in
the proof.  No assertion that local Weil pairings settle (A.17.2), and
no counterexample to that assertion for actual noncuspidal sections,
has been established in this work.

\begin{statement}{A.18}{Refuted}{The local rigidity claim L16}
At a good ordinary prime an integral Tate representation, even with
its Weil pairing, need not determine the Legendre parameter or its
$j$-invariant.
\end{statement}
\begin{proof}
Over $\Q_{17}$ take $E_4$ and $E_9$.  Both discriminants are units;
both reductions have 16 points and trace 2, so are ordinary.  The map
\[
 X=x-5+4/x,\qquad Y=y(1-4/x^2)
\]
is a degree-two isogeny
$E_4\to E':Y^2=X(X+1)(X+9)$, with kernel $\{O,(0,0)\}$.
Substitution proves the equation, and the equation for $x$ over $X$
is quadratic.  Hensel lifting $4\pmod {17}$ gives $i^2=-1$ in
$\Q_{17}$.  The map $(X,Y)\mapsto(-X,Y/i)$ is an isomorphism
$E'\to E_9$.  Their composite $\phi$ has degree two.
Thus $T_{17}(\phi)$ is an integral Galois isomorphism with Weil
multiplier two.  Hensel lifting $6\pmod {17}$ gives
$h\in\Z_{17}^\times$ with $h^2=2$, so
$h^{-1}T_{17}(\phi)$ is symplectic.  However
\[
                j(E_4)=35152/9,\qquad j(E_9)=1556068/81.
\]
These are unequal.  The displayed morphisms prove precisely the claimed
failure; the symbolic identities and finite counts are retained in
\texttt{code/legendre\_checks.py} and its output.
\end{proof}

\subsection{The modular unit, central lifting, and the Fermat twist}

\begin{statement}{A.19}{Proved}
{The exact modular-unit argument}
Put $U'=U\setminus\{-1,\tfrac12,2,t^2-t+1=0\}$.  Its unit
\[
 u(t)=1-1728/j(t)=
       \frac{(t-2)^2(t+1)^2(2t-1)^2}{4(t^2-t+1)^3}
\tag{A.19.1}
\]
has value one at each old cusp $0,1,\infty$.
A noncuspidal Selmer section of $U$ admitting a Selmer lift to $U'$
which is integral on $U'$ at almost every prime and whose local cusp
centres all belong to the old boundary is a rational point section.
The unrestricted lifting equivalence and the phrase ``whatever the
cusp labels'' in L17 have not been proved.
\end{statement}
\begin{proof}
Equation (A.19.1) gives zeros of order two at $-1,1/2,2$ and poles of
order three at the two roots of $t^2-t+1$.  The latter form one
degree-two closed point over $\Q$.  Pull back its Kummer class along
the lifted section.  At an interior local point the value is the
actual nonzero number $u(x_p)$.  At an old cusp, regularity and
nonvanishing make the Kummer class its residue value one, independently
of the tangent parameter.  Integrality on $U'$ makes these local values
units almost everywhere.  By A.10 the global completed class is one
$a\in\Q^\times$.  A.2 then puts all finite local centres in the
finite fibre $u^{-1}(a)$ on the completion.  Finite-support recognition
gives a point or cusp section upstairs.  A cusp in this fibre belongs
to the old boundary, since the new boundary consists of the zeros and
poles of $u$.  It would push to a cuspidal section of $U$, contrary to
the stated noncuspidality.  Hence the section is a rational point.

The two specific missing assertions in L17 are now visible in the
maps.  Integrality tested by $t,1-t$ on $U$ does not also test distance
from the five new punctures.  Moreover, at a new cusp the pullback of
$u$ records a tangent Kummer class with a nonzero zero or pole order;
it is not its nonzero residue value, so the finite-fibre inference used
above is unavailable.  These observations locate gaps; they do not
construct a nonpoint section satisfying L17.

The source AP-D.27 proves the reciprocity criterion for a completed
Kummer class and the associated \emph{cuspidally central} lift.  It
requires the restricted-product condition at every finite modulus in
addition to annihilation by all dual characters.  Its final sentence
leaves the full nonabelian lifting problem unresolved
\cite{LegAP2026}.  The comparison map is the quotient of the kernel
of $\pi_1(U')\to\pi_1(U)$ by the appropriate commutators to its
central Tate-module quotient.  Solving the embedding problem in that
quotient does not supply a section of $\pi_1(U')$ by any map constructed
here.  At the degree-two puncture the coefficient is
$\operatorname{Res}_{\Q(\sqrt{-3})/\Q}\mathbb G_m$ with its
restriction and corestriction maps; it is not two separately rational
punctures.
\end{proof}

\begin{statement}{A.20}{Refuted}
{The quadratic twist in the Fermat conductor assertion}
The raw residual Legendre representation of a Selmer section of the
Fermat quintic need not be unramified outside $2\cdot5$.
LM6-E.8's assertion after the specified quadratic twist is not refuted.
\end{statement}
\begin{proof}
On $X^5+Y^5+Z^5=0$ use the rational cusp $[1:-1:0]$ and parameter
$v=Z/X$.  The map to the tripod is
$t=-(X/Z)^5=-v^{-5}$ and $1-t=-(Y/Z)^5$.
Take the global cuspidal section with rational tangent $v=7$.
The infinity monodromy of the Legendre family is
\[
 C=\begin{pmatrix}1&-2\\2&-3\end{pmatrix}=-I+M,
          \qquad M^2=0.
\]
The cusp index of the cover is five, so its mod-five inertia is
$C^5=-I\pmod5$.  The tangent Kummer cocycle of 7 has valuation one
at 7, and thus pulls this nontrivial scalar inertia into the local
section image.  The tangent-one reference is unramified at 7: the
marked model and its constants have good reduction there.  The raw
residual representation is consequently ramified at 7.  Its quadratic
tangent character cancels this scalar, exactly the twist retained in
LM6-E.8 \cite{LegLM62026}.  Any use of L6's conductor claim must retain
that character or explicitly include its ramified primes in the set
being bounded.
\end{proof}

\subsection{Rank-zero twists and an explicit positive-rank exclusion}

\begin{statement}{A.21}{Proved}
{An honest point in an analytic-rank-zero cubic class}
The point $t=1/2$ lifts to $(1:1:1)$ on
$C_2:X^3+Y^3=2Z^3$.  This genus-one curve has Jacobian of
Mordell--Weil rank zero and analytic rank zero.  Thus honest points
need not belong to positive-rank cubic classes.
\end{statement}
\begin{proof}
With origin $(1:-1:0)$, the mutually inverse rational maps
\[
 x=\frac{6Z}{X+Y},\quad y=\frac{9(X-Y)}{X+Y},
 \qquad (X:Y:Z)=(9+y:9-y:3x)
\]
identify $C_2$ with $E:y^2=x^3-27$; substitution gives the equations,
and the maps extend over the smooth completions.  The displayed point
goes to $(3,0)$.

Here is an elementary rank calculation.  Translate $x=X+3$ to obtain
$E_0:y^2=X(X^2+9X+27)$, with 2-isogenous partner
$E_1:y^2=X(X^2-18X-27)$.  For $E_{a,b}:y^2=X(X^2+aX+b)$ define
$\alpha(O)=1$, $\alpha((0,0))=b$, and $\alpha((X,y))=X$ in
$\Q^\times/\Q^{\times2}$.  The divisor of $X$ is twice the
2-torsion divisor, so the chord law makes $\alpha$ a homomorphism.
For the isogeny
\[
 (X,y)\longmapsto
       (X+a+b/X,\ y(1-b/X^2))
\]
and its dual, the latter's $X$-coordinate is $(y'/(2X'))^2$.
Thus the dual image lies in $\ker\alpha$.  Conversely, if $X=r^2$,
the possible preimage coordinates solve
$X'^2-(2a+4X)X'+(a^2-4b)=0$, whose discriminant is
$(4y/r)^2$; taking $y'=\pm2rX'$ gives the desired preimage.
This proves the kernel assertion, including the special points by
continuation of the isogeny.

Valuations in the integral equation show that the image squareclasses
have squarefree representatives dividing $b$.  For $E_0$, the
quadratic $X^2+9X+27$ is everywhere positive on $\R$, so $X\ge0$.
The image is therefore $\{1,3\}$, with 3 supplied by $(0,0)$.
For $E_1$ the image contains $\{1,-3\}$ and is contained in
$\{1,-1,3,-3\}$.  The class 3 would give coprime integers $U,V$
and an integer $N$ satisfying
\[
                 N^2=3U^4-18U^2V^2-9V^4.
\]
Reduction modulo 3 forces $3\mid N,U$.  Write $N=3n,U=3u$ and
divide by 9.  Then
$n^2=27u^4-18u^2V^2-V^4$, which modulo 3 forces $3\mid V$,
a contradiction.  Since the image is a subgroup containing $-3$,
this also excludes $-1$, and its image is exactly $\{1,-3\}$.
The nonzero point in the dual kernel has nontrivial class $-3$.
The two kernel/cokernel sequences therefore give
\[
 |E_0(\Q)/2E_0(\Q)|=2\cdot2/2=2.
\]
Its rational 2-torsion has order two, so Mordell--Weil gives rank zero.

For analytic rank no numerical approximation is required.  The
associated form is $f(\tau)=\eta(6\tau)^4$, as identified explicitly
in \cite[Introduction and the $N=1$ case of the main theorem]{LegKanekoSakai2013}.
For $y>0$ all factors in
$\eta(6iy)^4=e^{-2\pi y}\prod_{n\ge1}(1-e^{-12\pi ny})^4$
are positive.  Its decay at infinity and the eta transformation at
zero make the Mellin integral convergent, and
\[
                     L(E,1)=2\pi\int_0^\infty\eta(6iy)^4\,dy>0.
\]
This proves analytic rank zero.
\end{proof}

\begin{statement}{A.22}{Proved}{The rank-zero part of A10}
Let $V\to U$ be a geometrically connected finite etale cover whose
completion admits a specified nonconstant $\Q$-morphism to an elliptic
curve of analytic rank zero.  Every Selmer section lifting to $V$ is
a point or cuspidal section.  This applies to a diagonal cubic when
its Jacobian has analytic rank zero, and to the explicit quartic and
sextic elliptic quotients below when their analytic rank is zero.
\end{statement}
\begin{proof}
An elliptic curve of analytic rank zero has finite Mordell--Weil and
Tate--Shafarevich groups; the finite-descent consequence used here is
exactly Stoll's Corollary 6.3.  Stoll's Proposition 8.5 and Theorem 8.6
apply to a nonconstant map to it \cite{LegStoll2007,LegStollErrata2017}.
The local centres of a lifted section map to one rational point of
that elliptic curve.  The inverse image is a finite subscheme of the
completion.  Finite-support recognition therefore gives an interior
point section or a boundary section of $V$, and its pushforward has
the asserted form on $U$.

We spell out maps so that the different Fermat exponents are not
identified without comparison.  For the cubic
$aX^3+bY^3=cZ^3$, let $D$ be its rational hyperplane divisor of
degree three.  The morphism
$P\mapsto\mathcal O(3P-D)$ to its Jacobian is defined over $\Q$;
after any geometric origin it is a translate of multiplication by
three and has degree nine.  For a quartic use affine coordinates
$au^4+bv^4=c$ and put $w=u^2$, $A=-b/a$, $B=c/a$.
Then $w^2=Av^4+B$, and
\[
       x=w^2/v^2,\qquad y=w(Av^4-B)/v^3
\]
give a nonconstant map to
$y^2=x^3-4ABx=x^3+(4bc/a^2)x$.
For the sextic $au^6+bv^6=c$, the map
\[
                  x=-(b/a)v^2,\qquad y=-(b/a)u^3
\]
lands on $y^2=x^3+b^2c/a^3$ and is nonconstant.
Direct substitution proves both equations; the rational maps extend
to the smooth proper completions.  These are the actual morphisms
through which an analytic-rank-zero entry of the permitted twist
table supplies the asserted finite support.
\end{proof}

\begin{statement}{A.23}{Proved}
{Primary profinite multiples that do not lift}
On $C_9:X^3+Y^3=9Z^3$ put $O=(1:-1:0)$,
$P=(1:2:1)$, and $V=C_9\setminus\{XYZ=0\}$.
The point $P$ has infinite order.  For every prime $\ell\ge5$ and
every $0\ne a\in\Z_\ell$, the element $a e_\ell P$ of
$\widehat{C_9(\Q)}$, where $e_\ell\in\widehat\Z$ is the
$\ell$-primary idempotent, is not the compactified Kummer class of
any Selmer section of $V$.
\end{statement}
\begin{proof}
The mutually inverse maps
\[
 x=\frac{12Z}{X+Y},\quad y=\frac{12(X-Y)}{X+Y},
 \qquad (X:Y:Z)=(12+y:12-y:2x)
\]
identify $(C_9,O)$ with $E:y^2=x^3-48$ and take $P$ to $(4,-4)$.
The finite etale map to the tripod is
$t=X^3/(9Z^3)$, $1-t=Y^3/(9Z^3)$, of degree nine, with mod-cube
Kummer class $(1/9,1/9)$.  The exact counts
$\#E(\mathbb F_5)=6$ and $\#E(\mathbb F_{13})=21$ bound its
rational torsion order by three.  Reduction is injective on rational
torsion at these primes: the kernel's formal logarithm identifies it
with a subgroup of the torsion-free additive group $p\Z_p$.
The chord law gives $2P=(28,148)\ne-P$, so $P$ has infinite order.
Mordell--Weil now makes $a e_\ell P$ nonzero in the profinite completion.

Take good primes $p\equiv2\pmod3$ and $p\equiv1\pmod\ell$.
They have positive density.  Cubing permutes $\mathbb F_p$, so the
quadratic-character sum for $x^3-48$ is zero and
$\#E(\mathbb F_p)=p+1\equiv2\pmod\ell$.
The group $E(\Q_p)$ is an extension of this finite group by a pro-$p$
group.  It has no $\ell$-primary factor, hence the localization of
$a e_\ell P$ is zero.  Its full local Kummer map is injective:
$E(\Q_p)$ is compact with an open subgroup isomorphic to $\Z_p$,
so $\bigcap_n nE(\Q_p)=0$.  A Selmer section of $V$ with the proposed
compactified class would therefore have local centre $O$ at all these
primes, and thus be locally cuspidal there.

Push it along $V\to U$.  If the image were globally noncuspidal,
AP-D.25's density-one integral-point set would contradict this
positive-density cusp set.  Thus the image is globally cuspidal.
For the actual transfer back, identify $\pi_1(V)$ with the open
subgroup $H_V\subset\pi_1(U)$.  After the relevant conjugation, its
section image lies in $H_V\cap D_c$, the decomposition group of an
extension of the cusp valuation in the normalization $C_9$.
Surjectivity onto $G_\Q$ forces the residue field of that boundary
point to be $\Q$.  The only rational deleted point of $C_9$ is $O$:
$X=0$ or $Y=0$ would make 9 a rational cube, and $Z=0$ gives only
$(1:-1:0)$ over $\Q$.  The section is therefore cuspidal at $O$.
Compactification kills its inertia, so its compactified Kummer class
is zero, contradicting $a e_\ell P\ne0$.
\end{proof}
This determines an explicit uncountable family of excluded profinite
multiples in a positive-rank class.  It does not classify all profinite
multiples, particularly those with infinitely many nonzero primary
components.  No lift is inferred merely from a Mordell--Weil
completion element.

\begin{statement}{A.24}{Open}{The remaining Legendre route to A1}
This route has not proved cofinite integrality for every noncuspidal
section, the vanishing of the multiplier class (A.17.2) for every
section in A.14, or recognition of a full section from its image under
the Legendre representation.  Neither assertion of A1 follows from
the results above.  The unclassified profinite multiples in A10 and
the full nonabelian lifting assertion in A8 also remain unresolved.
\end{statement}
The maps at these remaining steps are specified: reduction of the
integral representations in A.14, determinant on the torsor of
intertwiners in A.17, and the full continuous homomorphism
$\Pi\to H$ whose kernel is the relevant congruence kernel.  In
particular, A.18 prevents replacing the last step by the local
Tate-module assertion L16.  The precise comparisons proved here do
not yield an inverse to any of those quotient maps.
